\section{Introduction}
\label{sec:introduction}

Inductive and coinductive types are central to both functional programming and type theory. Semantically, an inductive type is the initial algebra of a functor, and a coinductive type is the terminal coalgebra of a functor. However, not every functor has such fixed-points: the functor \AgdaBound{X}\;\AgdaSymbol{↦}\;\AgdaSymbol{(}\AgdaBound{X}\;\AgdaSymbol{→}\;\AgdaDatatype{⊥}\AgdaSymbol{)} has none, since any fixed-point would yield a proof of \AgdaDatatype{⊥}. A well-behaved class of functors is given by \emph{containers}~\cite{AbbottAG03,AbbottAG05}, which represent exactly the strictly positive functors on types. Every container has an initial algebra, the \emph{W-type}~\cite{MartinLof84,GambinoH03}, and a terminal coalgebra, the \emph{M-type}~\cite{AhrensCS15}. Containers have since been extended in several directions, for instance to indexed containers~\cite{AltenkirchGHMM15} and to derivatives of datatypes~\cite{AbbottAGM05}.

Many datatypes of interest, however, are not fixed-points of functors on types. \emph{Nested datatypes}~\cite{BirdM98} are datatypes whose recursive occurrences are applied to arguments other than the type parameter being defined. A standard example is \AgdaDatatype{Bush}\;\AgdaBound{A}\;\AgdaSymbol{≃}\;\AgdaRecord{⊤}\;\AgdaDatatype{⊎}\;\AgdaBound{A}\;\AgdaFunction{×}\;\AgdaDatatype{Bush}\;\AgdaSymbol{(}\AgdaDatatype{Bush}\;\AgdaBound{A}\AgdaSymbol{)}. Since \AgdaDatatype{Bush} is applied to itself, it cannot be described as the initial algebra of a functor on types. Instead, it is the initial algebra of a functor on (strictly positive) functors: \AgdaBound{F}\;\AgdaSymbol{↦}\;\AgdaSymbol{λ}\AgdaBound{A}.\;\AgdaRecord{⊤}\;\AgdaDatatype{⊎}\;\AgdaBound{A}\;\AgdaFunction{×}\;\AgdaBound{F}\;\AgdaSymbol{(}\AgdaBound{F}\;\AgdaBound{A}\AgdaSymbol{)}. Nested datatypes have been used for several purposes: for instance, to represent first-order function types as terminal coalgebras~\cite{Altenkirch01}, and to enforce structural invariants such as well-scopedness of λ-terms~\cite{BirdP99b}. Their folds and unfolds have been studied extensively~\cite{BirdP99a,MatthesAU05,JohannG07}. More generally, datatypes of \emph{higher kinds} arise naturally in programming: monad transformers, free monads, and even the W-type itself are operations on containers. Such datatypes are the subject of recent work on the syntax and semantics of higher-kinded datatypes~\cite{JohannP19}, and have their roots in the type operators of System~Fω~\cite{Girard72}.

In this paper, we contruct a notion of container at higher kinds, and discuss if they have initial algebras and terminal coalgebras. Our starting point is the observation that, if we restrict the function variable of a signature of \AgdaDatatype{Bush} to be represented by containers, the signature becomes an endofunctor on the category of containers. We call the syntax of such endofunctors the \emph{second-order containers}, and we construct their fixed-points \emph{second-order W-types and M-types}. This picture extends to all finite kinds, leading us to \emph{higher-order containers}. (These should not be confused with the higher-order containers of Altenkirch, Levy and Staton~\cite{AltenkirchLS10}, which are exponential objects in the category of containers.) We find, however, that in the present setting fixed-points do not exist in general beyond the second order.

More concretely, our contributions are as follows.
\begin{itemize}
\item We give an explicit container representation of \AgdaDatatype{Bush}, obtained by solving a fixed-point equation in the category of containers. Its shapes and positions are defined mutually by induction-induction~\cite{NordvallForsbergS10} (Section~\ref{sec:bush}).

\item We introduce \emph{second-order containers}, which represent a class of strictly positive second-order polynomial functors, together with their    interpretations as functions on types and as endofunctors on containers (Section~\ref{sec:second}).

\item We construct, for every second-order container, an initial algebra \AgdaFunction{2W} by induction-induction and a terminal coalgebra \AgdaFunction{2M} by coinduction-induction (Section~\ref{sec:second}). \item We define \emph{higher-order containers} at every kind, as normal forms of a simply-typed λ-calculus with infinite products and coproducts, and we give both a type-level and a categorical semantics for them, as in the second-order case (Section~\ref{sec:higher}).

\item We show that, in this setting, fixed-points cannot exist beyond the second order: any higher-order fixed-point operator that allows us to construct an internal W-type would yield fixed-points for every type operator, including non-positive ones (Section~\ref{sec:nogo}).
  
\end{itemize}
